% GENERATED FILE. Edit canonical lessons, then run npm run generate:books.
% canonical-source-hash: fnv1a64:d3dd4ed86beebe9b
% canonical-lessons: ZH-C387-shouru, ZH-C387-cunkuan, ZH-C387-zhangdan, ZH-C387-wanggou, ZH-C387-yusuan

\chapter{Money In, Money Out}
\label{ch:zh-money-in-money-out}
\hlchaptermodality{387}

\begin{chapteropening}
\textbf{By the end of this chapter:} \emph{I can talk about income, savings, bills, shopping online and a budget.}

Complete the last lesson of chapter 387: I can talk about income, savings, bills, shopping online and a budget.
\end{chapteropening}

\section[shōurù]{shōurù — income}

\label{lesson:ZH-C387-shouru}



\begin{quote}
Before the new one: say the Mandarin for good value for money, then the Mandarin for to turn bad (of the weather).
\end{quote}

\subsection*{You'll want to know: shōurù}
\emph{shōurù} — \textquotedblleft{}income\textquotedblright{}.

\begin{quote}
\emph{yīnwèi tā de shōurù bù gāo suǒyǐ tā měitiān zài jiā zuòfàn} — His income isn't high, so he cooks at home every day.
\end{quote}

Its characters come later; for now it is a word for the ear.

\begin{cousinweb}
\emph{shōurù} is built from \emph{shōu} 'to gather' and \emph{rù} 'to go in', both inherited from Middle Chinese, \emph{syuw} and \emph{nyip} (Guangyun, in Baxter's transcription), which glosses the second 'to take in'.
\end{cousinweb}

\begin{grammarlens}[title={What you've built}]
One more word for giving reasons.
\end{grammarlens}

\subsection*{Your turn}
\begin{itemize}
\raggedright
  \item \emph{Say it:} \emph{shōurù}
  \item \emph{Say it:} \emph{shōurù}, once more
  \item \emph{Say it:} \emph{yīnwèi tā de shōurù bù gāo suǒyǐ tā měitiān zài jiā zuòfàn} — His income isn't high, so he cooks at home every day.
  \item \emph{Recall:} say the Mandarin for good value for money, then the Mandarin for to turn bad (of the weather), then say \emph{shōurù} again
\end{itemize}

\begin{tcolorbox}[breakable,colback=teal!4,colframe=teal!35!black,title={Before you move on}]
What does shōurù mean? (\textquotedblleft{}Income\textquotedblright{}.) Say it once more.
\end{tcolorbox}
\section[cúnkuǎn]{cúnkuǎn — savings (in the bank)}

\label{lesson:ZH-C387-cunkuan}



\begin{quote}
Before the new one: say the Mandarin for to turn bad (of the weather), then the Mandarin for income.
\end{quote}

\subsection*{You'll want to know: cúnkuǎn}
\emph{cúnkuǎn} — \textquotedblleft{}savings (in the bank)\textquotedblright{}.

\begin{quote}
\emph{wǒ yǒu yìdiǎnr cúnkuǎn} — I've got a little money saved.
\end{quote}

Its characters come later; for now it is a word for the ear.

\begin{cousinweb}
\emph{cúnkuǎn} is built from \emph{cún} 'to keep' and \emph{kuǎn} 'a sum', both inherited from Middle Chinese, \emph{dzwon} and \emph{khwanX} (Guangyun, in Baxter's transcription), which glosses the first 'to be present' and the second 'sincere'. The sense 'a sum of money' is not given there.
\end{cousinweb}

\begin{grammarlens}[title={What you've built}]
One more word for giving reasons.
\end{grammarlens}

\subsection*{Your turn}
\begin{itemize}
\raggedright
  \item \emph{Say it:} \emph{cúnkuǎn}
  \item \emph{Say it:} \emph{cúnkuǎn}, once more
  \item \emph{Say it:} \emph{wǒ yǒu yìdiǎnr cúnkuǎn} — I've got a little money saved.
  \item \emph{Recall:} say the Mandarin for to turn bad (of the weather), then the Mandarin for income, then say \emph{cúnkuǎn} again
\end{itemize}

\begin{tcolorbox}[breakable,colback=teal!4,colframe=teal!35!black,title={Before you move on}]
What does cúnkuǎn mean? (\textquotedblleft{}Savings (in the bank)\textquotedblright{}.) Say it once more.
\end{tcolorbox}
\section[zhàngdān]{zhàngdān — bill (to pay)}

\label{lesson:ZH-C387-zhangdan}



\begin{quote}
Before the new one: say the Mandarin for income, then the Mandarin for savings (in the bank).
\end{quote}

\subsection*{You'll want to know: zhàngdān}
\emph{zhàngdān} — \textquotedblleft{}bill (to pay)\textquotedblright{}.

\begin{quote}
\emph{zhège yuè de zhàngdān wǒ hái méi fù} — I haven't paid this month's bills yet.
\end{quote}

Its characters come later; for now it is a word for the ear.

\begin{cousinweb}
\emph{zhàngdān} is built from \emph{zhàng} 'an account' and \emph{dān} 'a single sheet'. The second is inherited from Middle Chinese \emph{tan} (Guangyun, in Baxter's transcription), glossed 'single'; the dictionary has no entry for the first, so its history is not settled.
\end{cousinweb}

\begin{grammarlens}[title={What you've built}]
One more word for giving reasons.
\end{grammarlens}

\subsection*{Your turn}
\begin{itemize}
\raggedright
  \item \emph{Say it:} \emph{zhàngdān}
  \item \emph{Say it:} \emph{zhàngdān}, once more
  \item \emph{Say it:} \emph{zhège yuè de zhàngdān wǒ hái méi fù} — I haven't paid this month's bills yet.
  \item \emph{Recall:} say the Mandarin for income, then the Mandarin for savings (in the bank), then say \emph{zhàngdān} again
\end{itemize}

\begin{tcolorbox}[breakable,colback=teal!4,colframe=teal!35!black,title={Before you move on}]
What does zhàngdān mean? (\textquotedblleft{}Bill (to pay)\textquotedblright{}.) Say it once more.
\end{tcolorbox}
\section[wǎnggòu]{wǎnggòu — to shop online}

\label{lesson:ZH-C387-wanggou}



\begin{quote}
Before the new one: say the Mandarin for savings (in the bank), then the Mandarin for a bill (to pay).
\end{quote}

\subsection*{You'll want to know: wǎnggòu}
\emph{wǎnggòu} — \textquotedblleft{}to shop online\textquotedblright{}.

\begin{quote}
\emph{wǒ juéde wǎnggòu bǐ qù shāngdiàn piányi} — I think shopping online is cheaper than going to the shops.
\end{quote}

Its characters come later; for now it is a word for the ear.

\begin{cousinweb}
\emph{wǎnggòu} is built from \emph{wǎng} 'a net' and \emph{gòu} 'to buy', both inherited from Middle Chinese, \emph{mjangX} and \emph{kuwH} (Guangyun, in Baxter's transcription), which glosses the second 'to buy back'. The word is a modern short form.
\end{cousinweb}

\begin{grammarlens}[title={What you've built}]
One more word for giving reasons.
\end{grammarlens}

\subsection*{Your turn}
\begin{itemize}
\raggedright
  \item \emph{Say it:} \emph{wǎnggòu}
  \item \emph{Say it:} \emph{wǎnggòu}, once more
  \item \emph{Say it:} \emph{wǒ juéde wǎnggòu bǐ qù shāngdiàn piányi} — I think shopping online is cheaper than going to the shops.
  \item \emph{Recall:} say the Mandarin for savings (in the bank), then the Mandarin for a bill (to pay), then say \emph{wǎnggòu} again
\end{itemize}

\begin{tcolorbox}[breakable,colback=teal!4,colframe=teal!35!black,title={Before you move on}]
What does wǎnggòu mean? (\textquotedblleft{}To shop online\textquotedblright{}.) Say it once more.
\end{tcolorbox}
\section[yùsuàn]{yùsuàn — budget}

\label{lesson:ZH-C387-yusuan}



\begin{quote}
Before the new one: say the Mandarin for a bill (to pay), then the Mandarin for to shop online.
\end{quote}

\subsection*{You'll want to know: yùsuàn}
\emph{yùsuàn} — \textquotedblleft{}budget\textquotedblright{}.

\begin{quote}
\emph{wǒmen de yùsuàn bú gòu suǒyǐ zuìhǎo mǎi èrshǒu de} — Our budget isn't enough, so we'd better buy second-hand.
\end{quote}

Its characters come later; for now it is a word for the ear.

\begin{cousinweb}
\emph{yùsuàn} is built from \emph{yù} 'beforehand' and \emph{suàn} 'to count', both inherited from Middle Chinese, \emph{yoH} and \emph{swanX} (Guangyun, in Baxter's transcription), which glosses the first, among other senses, 'before'.
\end{cousinweb}

\begin{grammarlens}[title={What you've built}]
One more word for giving reasons.
\end{grammarlens}

\subsection*{Your turn}
\begin{itemize}
\raggedright
  \item \emph{Say it:} \emph{yùsuàn}
  \item \emph{Say it:} \emph{yùsuàn}, once more
  \item \emph{Say it:} \emph{wǒmen de yùsuàn bú gòu suǒyǐ zuìhǎo mǎi èrshǒu de} — Our budget isn't enough, so we'd better buy second-hand.
  \item \emph{Recall:} say the Mandarin for a bill (to pay), then the Mandarin for to shop online, then say \emph{yùsuàn} again
\end{itemize}

\begin{tcolorbox}[breakable,colback=teal!4,colframe=teal!35!black,title={Before you move on}]
What does yùsuàn mean? (\textquotedblleft{}Budget\textquotedblright{}.) Say it once more.
\end{tcolorbox}
